%% kapa: exam=1 round=30 year=2019 period=1 form=A subject=civil_law subject_ko=민법 seq=16 number=16
%% source: Q-Net 감정평가사 기출문제 (cst00309) artlSeq=5205205 file=extracted/1008/5205205/2019년 제30회 감정평가사 1차시험 문제지/2019년 제30회 감정평가사 1차시험 문제지/2019년도 제30회 감정평가사 1차 1교시 A형.hwp
%% conversion: kordoc 4.24.51 -> tools/gen1.py (auto); review=proofread-sonnet; flags=none
\begin{question}{16}
법률행위의 취소에 관한 설명으로 옳지 않은 것은?
\begin{choices}
  \item 착오로 인하여 취소할 수 있는 법률행위를 한 자의 포괄승계인은 그 법률행위를 취소할 수 있다.
  \item 미성년자가 동의없이 단독으로 행한 법률행위를 그 법정대리인이 추인하는 경우, 그 추인은 취소의 원인이 소멸한 후에 하여야만 효력이 있다.
  \item 제한능력자가 제한능력을 이유로 법률행위를 취소한 경우, 그 행위로 인하여 받은 이익이 현존하는 한도에서 상환할 책임이 있다.
  \item 취소할 수 있는 법률행위를 추인한 후에는 이를 다시 취소하지 못한다.
  \item 취소권은 추인할 수 있는 날로부터 3년 내에, 법률행위를 한 날로부터 10년 내에 행사하여야 한다.
\end{choices}
\end{question}
\answer{2}
